\clearpage\NoSourceHead\SourceLeaf{405}
\SourceGroup{IV}
\HistoricTitle{[DARSTELLUNG VON DISKONTINUIERLICHEN FUNKTIONEN.]}
\HistoricCitation{[Aus Handbuch 19, Be, Kleine Aufsätze aus verschiedenen Theilen der Mathematik, Angefangen im May 1809, S. 277.]}
\noindent\hspace*{\parindent}Die discontinuirliche Function von \(t\), welche für \(t\) von 0 bis \(180^0\) mit \(\sin t\) übereinstimmt, von \(180^0\) bis \(360^0\) aber verschwindet, ist
\[
=\frac1\pi+\frac12\sin t-\frac2{1{.}3{.}\pi}\cos2t-\frac2{3{.}5{.}\pi}\cos4t-\frac2{5{.}7{.}\pi}\cos6t-\ \text{u.s.w.}
\]
\noindent oder, was dasselbe ist, wenn der Werth von \(90^0\) bis \(270^0\) verschwinden, von 0 bis \(90^0\) und von \(270^0\) bis \(360^0\) aber mit \(\cos t\) übereinstimmen soll, ist die Formel
\[
=\frac1\pi+\frac12\cos t+\frac2{1{.}3{.}\pi}\cos2t-\frac2{3{.}5{.}\pi}\cos4t+\frac2{5{.}7{.}\pi}\cos6t-\mathord{\cdot}\,\mathord{\cdot}\,\mathord{\cdot}.
\]
\noindent Man kann dies auch so ausdrücken:
\[
\frac2\pi-\frac4{3\pi}\cos u-\frac4{15\pi}\cos2u-\frac4{35\pi}\cos3u-\ \text{u.s.w.}
\]
\noindent gibt immer den positiven Werth von \(\sin\frac12u\).
\ClosingRule
\par\bigskip\noindent\hspace*{\parindent}Die Function von \(u\), welche von \(u=-A\) bis \(u=+A\) constant \(=1\) sein, für alle übrigen Werthe aber verschwinden soll, ist
\[
\frac A\pi+\frac{2\sin A}\pi\cdot\cos u+\frac{2\sin2A}{2\pi}\cdot\cos2u+\frac{2\sin3A}{3\pi}\cdot\cos3u+\ \text{u.s.w.}
\]
\ClosingRule
\clearpage\SourceLeaf{406}\ArithmeticHead{399}{R}{DARSTELLUNG VON DISKONTINUIERLICHEN FUNKTIONEN.}
\noindent\hspace*{\parindent}Die Function soll von \(u=0\) bis \(u=A\) durch \(A-u\), von \(u=-A\) bis \(u=0\) durch \(A+u\) dargestellt werden, ausserhalb dieser Grenzen aber verschwinden:
\[
\begin{aligned}
\frac1\pi\Bigl(\frac12AA+2(1-\cos A)\cos u+\frac24(1-\cos2A)\cos2u\qquad\\
+\frac29(1-\cos3A)\cos3u+\text{u.s.w.}\Bigr)
\end{aligned}
\]
\ClosingRule
\par\bigskip\noindent\hspace*{\parindent}Die Function soll constant sein \(=B-A\) für \(u=-A\) bis \(u=+A\), ferner von \(u=A\) bis \(u=B\) gleichförmig abnehmen, bis sie für \(u=B\) verschwindet; eben so soll sie wie \(u\) von \(-A\) zu \(-B\) übergeht gleichförmig von \(B-A\) zu 0 übergehen und ausserhalb der Grenzen \(-B\) und \(+B\), d. i. von \(-180^0\) bis \(-B\) und von \(+B\) bis \(+180^0\) verschwinden:
\[
\begin{aligned}
\frac1\pi\Bigl\{\frac12(BB-AA)+2(\cos A-\cos B)\cos u\qquad\\
+\frac24(\cos2A-\cos2B)\cos2u+\frac29(\cos3A-\cos3B)\cos3u+\text{etc.}\Bigr\}.
\end{aligned}
\]
\begin{center}BEMERKUNG.\end{center}
{\small\noindent\hspace*{\parindent}Auf der vorhergehenden Seite 276 des Handbuchs endet die S. 273 beginnende Abhandlung über den Lehrsatz von \textsc{Lagrange}, die Werke VIII, S. 80—83 abgedruckt ist, mit der Datierung »1847, Mai 13.«\par}
\begin{flushright}\textsc{Schlesinger}.\end{flushright}
\ClosingRule
